%% toc: P15 160 160.9
% Proofs for the second-derivative reading of Durrett PTE5 Theorem 3.3.21 (p. 136) and the trigonometric derivation of Exercise 3.3.10 (p. 133).
\section{Second derivatives, sine products and the weak-law derivative criterion}
\label{sec:annex-p15-w02c-p15-013-second-derivatives-and-products}

\begin{annexitem}{annex:p15-w02c-p15-013-quotient-limit}{thm:p15-w02c-p15-013-quotient-limit}{Cauchy's mean value step}
\begin{proof}
Write \(D(h)=f(h)-2f(0)+f(-h)\). Since \(D(-h)=D(h)\), it suffices to treat
\(h\downarrow0\). Fix \(h\in(0,\delta)\) and define on \([0,h]\)
\[
 F(s)=f(s)+f(-s)-2f(0),\qquad
 H(s)=h^2F(s)-s^2F(h).
\]
Both \(s\) and \(-s\) lie in \((-\delta,\delta)\) for \(s\in[0,h]\), so \(F\)
and \(H\) are differentiable on \([0,h]\), with
\(F'(s)=f'(s)-f'(-s)\). Also \(F(0)=0\), so \(H(0)=0\), and
\(H(h)=h^2F(h)-h^2F(h)=0\). Rolle's theorem gives \(\xi\in(0,h)\) with
\[
 0=H'(\xi)=h^2\bigl(f'(\xi)-f'(-\xi)\bigr)-2\xi F(h).
\]
Since \(\xi>0\) and \(F(h)=D(h)\),
\[
 \frac{D(h)}{h^2}=\frac{f'(\xi)-f'(-\xi)}{2\xi}
 =\frac12\left(\frac{f'(\xi)-f'(0)}{\xi}
             +\frac{f'(-\xi)-f'(0)}{-\xi}\right).
\]
The point \(\xi=\xi(h)\) satisfies \(0<\xi<h\), so \(\xi\to0\) as
\(h\downarrow0\). Differentiability of \(f'\) at \(0\) makes each of the two
difference quotients in parentheses tend to \(f''(0)\). Hence
\(D(h)/h^2\to f''(0)\).
\end{proof}
\end{annexitem}

\begin{annexitem}{annex:p15-w02c-p15-013-second-derivative-moment}{thm:p15-w02c-p15-013-second-derivative-moment}{From the derivative to the quotient, then to the moment}
\begin{proof}
\emph{Item~1.} Write \(v=\operatorname{Im}\varphi\). By
\cref{thm:p15-characteristic-algebra}, \(\varphi(0)=1\) and
\(\varphi(-t)=\overline{\varphi(t)}\), so \(u(-t)=u(t)\) and \(v(-t)=-v(t)\).
The numerator of \(Q(h)\) is
\[
 \varphi(h)-2+\varphi(-h)=\bigl(u(h)-2u(0)+u(-h)\bigr)
     +i\bigl(v(h)+v(-h)\bigr),
\]
and \(v(h)+v(-h)=0\), while \(u(0)=1\). Thus
\(Q(h)=(u(h)-2u(0)+u(-h))/h^2\) is real. Componentwise differentiation
(\cref{def:p15-characteristic-complex-conventions}) makes \(u\) differentiable
on \((-\delta,\delta)\) with \(u'=\operatorname{Re}\varphi'\). The real parts
of the quotients \((\varphi'(h)-\varphi'(0))/h\) are
\((u'(h)-u'(0))/h\), so \(u'\) is differentiable at \(0\) with
\(u''(0)=\operatorname{Re}\varphi''(0)\).
\Cref{thm:p15-w02c-p15-013-quotient-limit} with \(f=u\) gives
\(Q(h)\to\operatorname{Re}\varphi''(0)\). This limit is finite, so
\(\limsup_{h\downarrow0}Q(h)=\operatorname{Re}\varphi''(0)>-\infty\).
Item~2 will show that \(\varphi''(0)\) is real, so the limit is
\(\varphi''(0)\).

\emph{Item~2.} The hypothesis of \cref{thm:p15-characteristic-second-moment}
holds by item~1. That theorem gives \(\mathbb EX^2<\infty\) and
\(Q(h)\to-\mathbb EX^2\). Since limits are unique,
\(\operatorname{Re}\varphi''(0)=-\mathbb EX^2\). Now apply
\cref{thm:p15-characteristic-moment-derivatives} with \(n=2\). It shows that
\(\varphi\) has continuous derivatives through order \(2\) on \(\mathbb R\)
and that \(\varphi^{(2)}(0)=i^2\mathbb EX^2=-\mathbb EX^2\). On
\((-\delta,\delta)\) the function \(\varphi'\) of that theorem is the same
function as the \(\varphi'\) of the hypothesis, since a derivative at a point
is unique. Hence its derivative at \(0\) is the \(\varphi''(0)\) of the
hypothesis, and \(\varphi''(0)=-\mathbb EX^2\) is real and at most \(0\).
Together with item~1, \(Q(h)\to\varphi''(0)\).

\emph{Item~3.} Since \(|X|\leq1+X^2\), \(\mathbb E|X|\leq1+\mathbb EX^2<\infty\).
\Cref{thm:p15-characteristic-moment-derivatives} with \(n=2\) gives
\(\varphi^{(j)}(t)=\mathbb E[(iX)^je^{itX}]\) for \(j=1,2\) and all real
\(t\), with continuous derivatives through order \(2\). At \(j=1\), \(t=0\)
this is \(\varphi'(0)=i\,\mathbb EX\). At \(j=2\) it is
\(\varphi''(t)=-\mathbb E[X^2e^{itX}]\).
\end{proof}
\end{annexitem}

\begin{annexitem}{annex:p15-w02c-p15-013-four-conditions}{thm:p15-w02c-p15-013-four-conditions}{The cycle of implications}
\begin{proof}
The cycle (D)\(\Rightarrow\)(S)\(\Rightarrow\)(U)\(\Rightarrow\)\(\mathbb EX^2<\infty\)\(\Rightarrow\)(D)
is proved.

(D)\(\Rightarrow\)(S): item~1 of
\cref{thm:p15-w02c-p15-013-second-derivative-moment} gives
\(Q(h)\to\varphi''(0)\in\mathbb R\).

(S)\(\Rightarrow\)(U): if \(Q(h)\to L\in\mathbb R\), the upper limit as
\(h\downarrow0\) equals \(L>-\infty\).

(U)\(\Rightarrow\mathbb EX^2<\infty\): this is the forward direction of
\cref{thm:p15-characteristic-second-moment}.

\(\mathbb EX^2<\infty\Rightarrow\)(D): by
\cref{thm:p15-characteristic-moment-derivatives} with \(n=2\), \(\varphi\) is
differentiable on \(\mathbb R\) and \(\varphi'\) is differentiable at every
point, in particular at \(0\). So (D) holds with any \(\delta>0\).

When the conditions hold, item~2 of
\cref{thm:p15-w02c-p15-013-second-derivative-moment} gives
\(\varphi''(0)=-\mathbb EX^2\), and item~1 gives
\(\lim_{h\to0}Q(h)=\varphi''(0)\).

For \(g(t)=|t|\) and \(h\neq0\),
\((g(h)-2g(0)+g(-h))/h^2=2|h|/h^2=2/|h|\to+\infty\), so the upper limit in
(U) is \(+\infty\). The one-sided quotients \((g(h)-g(0))/h\) equal \(1\) for
\(h>0\) and \(-1\) for \(h<0\), so \(g\) is not differentiable at \(0\), and
(D) fails. Hence the step (U)\(\Rightarrow\)(D) depends on the
characteristic-function structure, which enters through
\cref{thm:p15-characteristic-second-moment}.
\end{proof}
\end{annexitem}

\begin{annexitem}{annex:p15-w02c-p15-013-order-contrast}{thm:p15-w02c-p15-013-order-contrast}{Exponential inequalities at the origin}
\begin{proof}
Two inequalities for real \(y\geq0\) are used:
\[
 1-y\leq e^{-y}\leq1-y+\tfrac12y^2 .
\]
For the first, \(a(y)=e^{-y}-1+y\) has \(a(0)=0\) and
\(a'(y)=1-e^{-y}\geq0\). For the second, \(b(y)=1-y+\tfrac12y^2-e^{-y}\) has
\(b(0)=0\) and \(b'(y)=-1+y+e^{-y}=a(y)\geq0\).

Write \(\varphi=\varphi_\alpha\). For \(t\neq0\), the chain rule gives
\(\varphi'(t)=-\alpha|t|^{\alpha-1}\operatorname{sgn}(t)\,e^{-|t|^\alpha}\),
which is continuous on \(t\neq0\). At \(0\), the first inequality with
\(y=|h|^\alpha\) gives \(0\leq1-\varphi(h)\leq|h|^\alpha\), so
\(|(\varphi(h)-\varphi(0))/h|\leq|h|^{\alpha-1}\to0\), since \(\alpha>1\).
Thus \(\varphi'(0)=0\). Also \(|\varphi'(t)|\leq\alpha|t|^{\alpha-1}\to0\) as
\(t\to0\), so \(\varphi'\) is continuous at \(0\) as well.

Since \(\varphi\) is real and even,
\(Q(h)=2(e^{-|h|^\alpha}-1)/h^2\). For \(0<|h|\leq1\), put
\(y=|h|^\alpha\in(0,1]\). The second inequality gives
\(1-e^{-y}\geq y-\tfrac12y^2\geq\tfrac12y\), using \(y\leq1\). Hence
\[
 Q(h)\leq-\frac{2\cdot\tfrac12|h|^\alpha}{h^2}=-|h|^{\alpha-2},
\]
and \(|h|^{\alpha-2}\to\infty\) as \(h\to0\), because \(\alpha<2\). So
\(Q(h)\to-\infty\). If \(\varphi''(0)\) existed, item~1 of
\cref{thm:p15-w02c-p15-013-second-derivative-moment} would give a finite
limit of \(Q(h)\); so \(\varphi''(0)\) does not exist. If
\(\mathbb EX^2\) were finite, \cref{thm:p15-characteristic-second-moment}
would give \(Q(h)\to-\mathbb EX^2\), again a finite limit. Hence
\(\mathbb EX^2=\infty\).
\end{proof}
\end{annexitem}

\begin{annexitem}{annex:p15-w02c-p15-013-sine-product}{thm:p15-w02c-p15-013-sine-product}{Induction on the number of halvings}
\begin{proof}
\emph{The double-angle identity.} For real \(x\),
\(e^{2ix}=(e^{ix})^2=(\cos x+i\sin x)^2\). Comparing imaginary parts gives
\(\sin2x=2\sin x\cos x\).

\emph{The finite identity.} For \(n=1\), the double-angle identity at
\(x=t/2\) is \(\sin t=2\sin(t/2)\cos(t/2)=2^1\sin(t/2)P_1(t)\). Suppose the
identity holds for some \(n\geq1\). The double-angle identity at
\(x=t/2^{n+1}\) gives \(\sin(t/2^n)=2\sin(t/2^{n+1})\cos(t/2^{n+1})\).
Substituting,
\[
 \sin t=2^n\cdot2\sin(t/2^{n+1})\cos(t/2^{n+1})\,P_n(t)
       =2^{n+1}\sin(t/2^{n+1})\,P_{n+1}(t),
\]
since \(P_{n+1}(t)=P_n(t)\cos(t/2^{n+1})\). This proves the identity for all
\(n\geq1\).

\emph{The quotient form.} Let \(t\neq0\) and \(2^n>|t|/\pi\). Then
\(x_n=t/2^n\) satisfies \(0<|x_n|<\pi\), so \(\sin x_n\neq0\), as the zeros of
the sine are the integer multiples of \(\pi\). Division gives
\(P_n(t)=\sin t/(2^n\sin x_n)\).

\emph{The limit.} The imaginary part of \(e^{ix}-1-ix\) is \(\sin x-x\), so
\cref{thm:p15-characteristic-taylor-remainder} with \(n=1\) gives
\(|\sin x-x|\leq|e^{ix}-1-ix|\leq x^2/2\). For \(x\neq0\) this yields
\(|\sin x/x-1|\leq|x|/2\). For fixed \(t\neq0\), \(x_n=t/2^n\to0\), so
\(\sin x_n/x_n\to1\). For all large \(n\),
\[
 P_n(t)=\frac{\sin t}{t}\cdot\frac{x_n}{\sin x_n}\longrightarrow\frac{\sin t}{t}.
\]
At \(t=0\), every factor is \(\cos0=1\), so \(P_n(0)=1\) for every \(n\),
which is the value assigned to \((\sin t)/t\) at \(0\).

\emph{Failure of uniformity.} At \(t=2^n\pi\), the factor with index \(m\) is
\(\cos(2^{n-m}\pi)\). For \(m=n\) it is \(\cos\pi=-1\). For \(1\leq m<n\),
\(2^{n-m}\) is an even integer, so the factor is \(1\). Thus
\(P_n(2^n\pi)=-1\). Since \(\sin(2^n\pi)=0\) and \(2^n\pi\neq0\), the limit
function is \(0\) there. Hence
\(\sup_{t}|P_n(t)-(\sin t)/t|\geq|-1-0|=1\) for every \(n\), and the
convergence is not uniform on \(\mathbb R\).
\end{proof}
\end{annexitem}

\begin{annexitem}{annex:p15-w02c-p15-013-half-pi}{thm:p15-w02c-p15-013-half-pi}{Half angles in the first quadrant}
\begin{proof}
First, \(c_1=\sqrt2/2=\cos(\pi/4)=\cos(\pi/2^2)\). Suppose
\(c_m=\cos(\pi/2^{m+1})\). The angle \(x=\pi/2^{m+1}\) lies in
\((0,\pi/2]\), so \(x/2\in(0,\pi/4]\) and \(\cos(x/2)>0\). The identity
\(\cos x=2\cos^2(x/2)-1\), which is the real part of
\(e^{ix}=(e^{ix/2})^2\), gives \(\cos^2(x/2)=(1+\cos x)/2\). Taking the
positive root, \(\cos(\pi/2^{m+2})=\sqrt{(1+c_m)/2}=c_{m+1}\). By induction
\(c_m=\cos(\pi/2^{m+1})\) for all \(m\geq1\).

Apply \cref{thm:p15-w02c-p15-013-sine-product} at \(t=\pi/2\). The factor
with index \(m\) is \(\cos(\pi/2^{m+1})=c_m\), and
\((\sin(\pi/2))/(\pi/2)=2/\pi\). So \(\prod_{m\geq1}c_m=2/\pi\). The displayed
radicals follow: \(c_2=\sqrt{(1+\sqrt2/2)/2}=\sqrt{2+\sqrt2}/2\), and in
general \(c_m=r_m/2\) with \(r_1=\sqrt2\) and \(r_{m+1}=\sqrt{2+r_m}\), since
\(\sqrt{(1+r_m/2)/2}=\sqrt{2+r_m}/2\).
\end{proof}
\end{annexitem}

\begin{annexitem}{annex:p15-w02c-p15-013-binary-uniform}{thm:p15-w02c-p15-013-binary-uniform}{Product rule, trigonometry and uniqueness}
\begin{proof}
Take versions of the \(\varepsilon_j\) that take only the values \(\pm1\)
at every outcome. Changing each variable on a null set does not change
any joint law. Then \(|\varepsilon_j2^{-j}|=2^{-j}\), and
\(\sum_j2^{-j}=1\), so the series converges absolutely at every outcome.
Its partial sums \(V_n\) converge to \(V\) everywhere, hence almost surely.

The fair sign \(\varepsilon_j\) has characteristic function \(\cos t\)
(\cref{thm:p15-characteristic-elementary-examples}). By the scaling rule of
\cref{thm:p15-characteristic-algebra} with \(a=2^{-j}\), \(b=0\), the
variable \(\varepsilon_j2^{-j}\) has characteristic function
\(\cos(t/2^j)\). By the product rule for independent summands, \(V_n\) has
characteristic function \(P_n(t)\). Item~3 of
\cref{thm:p15-criteria-limit-operations} gives
\(\varphi_V(t)=\lim_nP_n(t)\) for every real \(t\). By
\cref{thm:p15-w02c-p15-013-sine-product}, this limit is \((\sin t)/t\), with
value \(1\) at \(0\). By \cref{thm:p15-characteristic-elementary-examples}
with \(c=1\), \((\sin t)/t\) is the characteristic function of the uniform
law on \((-1,1)\). By the uniqueness assertion of
\cref{thm:p15-characteristic-inversion-uniqueness}, two Borel probability
laws with the same characteristic function are equal. So \(V\) is uniform
on \((-1,1)\).
\end{proof}
\end{annexitem}
